\section{Bohr sets and sets of positive measure}
\label{sec:density}

\subsection{A regularity lemma}
The goal of this section is to prove \cref{prop:regularity_lemma}. As mentioned in the introduction, this argument has its genesis in Bourgain \cite{bourgain-roth}. Bourgain's ideas were elaborated by Tao \cite{tao}, who proved Roth's theorem in compact abelian groups that are 2-divisible; and by Bergelson-Host-McCutcheon-Parreau \cite[Theorem 4.1]{bhmp}, who proved Roth's theorem for dilations on the torus $\R / \Z$. We streamline and generalize Bergelson-Host-McCutcheon-Parreau's argument to deal with homomorphisms on arbitrary compact abelian groups. This generalization requires non-trivial modifications; especially, we will make use of \cref{lem:translate} and \cref{lem:kernel_a}. 

%\Hnote{Define $d(t), J(f)$, etc.}

% For $f \in L^\infty(G)$, we write
% \[
% J(f) = \iint_{G^2} f(x)f(x+ \phi(y))f(x+ \psi(y)) \, d\mu(x) d\mu(y)
% \]

\begin{lemma}[cf. {\cite[Lemma 4.2]{bhmp}}]
\label{lem:J}
Let $\phi, \psi: G \rightarrow G$ be continuous homomorphisms such that $\phi(G), \psi(G)$ and $(\phi-\psi)(G)$ have finite indexes in $G$. For $f \in L^\infty(G)$, define
\[
J(f) = \iint_{G^2} f(x)f(x+ \phi(y))f(x+ \psi(y)) \, d\mu(x) d\mu(y).
\]
Then for any measurable functions $f, g: G \to [0,1]$, 
    \begin{equation*} \label{eq:c}
        |J(f) - J(g)| \ll \| \widehat{f} - \widehat{g} \|_\infty,
    \end{equation*}
where the implicit constant depends only on the aforementioned indexes. 
%$[G: \phi(G)], [G: \psi(G)]$ and $[G: (\phi-\psi)(G)]$.		
\end{lemma}

\begin{proof}
We have
\begin{eqnarray*}
J(f) - J(g) &=& \iint_{G^2} (f-g)(x) \cdot f(x+ \phi(y)) \cdot f(x+ \psi(y)) \, d\mu(x) d\mu(y) \\
& & + \iint_{G^2} g(x) \cdot ( f-g) (x+ \phi(y)) \cdot f(x+ \psi(y)) \, d\mu(x) d\mu(y) \\
& & + \iint_{G^2} g(x) \cdot g(x+ \phi(y)) \cdot (f - g) (x + \psi(y)) \, d\mu(x) d\mu(y).
\end{eqnarray*}
The lemma now follows from Lemma \ref{lem:bourgain_lemma_2} and the assumptions on $\phi$ and $\psi$.
\end{proof}


% \textcolor{red}{Choose between $\lVert K \rVert_{A(G)}, \lVert \widehat{K} \rVert_1$ (and maybe $\lVert K \rVert_{\infty}$.)}

\begin{proposition}[Regularity Lemma]
\label{prop:regularity_lemma} Let $f: G \to [0, 1]$ be a measurable function with $\int_G f \, d \mu = \delta >0$. Let $\phi, \psi: G \to G$ be continuous homomorphisms such that $\phi(G), \psi(G)$ and $(\phi- \psi)(G)$ have finite indexes in $G$. Then for every $\epsilon > 0$, there exist a constant $C$ that depends only on $\delta, \epsilon$ and the indexes above, a kernel $K: G \to \R_{\geq 0}$, and a decomposition $f = f_{st} + f_{er} + f_{un}$ such that
\begin{enumerate}
    \item $\lVert K \rVert_{\infty} < C$,
    \item $\lVert f_{st} \rVert_{\infty} \leq 1$, $\lVert f_{er} \rVert_{\infty} \leq 2$ and $\lVert f_{un} \rVert_{\infty} \leq 2$,
    \item $\lVert f_{er} \rVert_2 < \epsilon$,
    \item $\lVert \widehat{f}_{un} \rVert_{\infty} \lVert \widehat{K} \rVert_1 < \epsilon$,
    \item $J'(f_{st}) := \displaystyle \iint_{G^2} f_{st}(x) f_{st}(x+ \phi(t)) f_{st}(x+ \psi(t)) K(t) \, d\mu(x) d\mu(t) \geq \delta^3 - \epsilon$.
\end{enumerate}

%\textcolor{red}{Is $\lVert \hat{k} \rVert_1 = \lVert k \rVert_{A(G)} =  \lVert k \rVert_{\infty}?$ This equality is claimed in \cref{lem:kernel_a}.}
\end{proposition}
\begin{proof}
For $t \in G$, let
		\[
			d(t) := \max \left( \| \widehat{f} - \widehat{f_{t}} \|_\infty, \| \widehat{f} - \widehat{f_{\phi(t)}} \|_\infty, \| \widehat{f} - \widehat{f_{\psi(t)}} \|_\infty  \right).
		\]
Fixing $\epsilon > 0$, we define sequences $\eta_n \in (0, 1/2]$, $\kappa_n \in (0, \infty)$ and finite sets $\Lambda_n \subseteq \Gamma$ recursively as follows: 

First set $\eta_0 = 1/2$. For $n \geq 0$, \cref{lem:translate} implies that there exists a set $\Lambda_n \in \Gamma$ with $|\Lambda_n| \leq 12/\eta_n^2$ such that $D(\eta_n) := \{t \in G: d(t) \leq \eta_n\}$ contains a Bohr set $B(\Lambda_n; \eta_n)$.  
For $\eta \in (0, 1/2]$, define $\nu(\eta) = (C_0 \eta)^{12/\eta^2}$ where $C_0$ is the constant found in \cref{lem:kernel_a}; in particular, $\nu(\eta_n) = (C_0 \eta_n)^{12/\eta_n^2} \leq (C_0 \eta_n)^{|\Lambda_n|}$. Put
\[
         \kappa_n = \nu(\eta_n)^{-1/2} \text{ and } \eta_{n+1} = \min \left \{\eta_n, \frac{\epsilon^2}{4 \kappa_n^2}, \epsilon \nu \left( \frac{\epsilon}{2 \kappa_n} \right) \right\}.
\]
% Fixing $\epsilon > 0$, we define decreasing sequences $\eta_n \in (0, 1/2]$, $\kappa_n \in (0, \infty)$ and increasing nested sequence $\Lambda_n \subseteq \Gamma$ as follows: First set $\eta_0 = 1/2$ and for each $n \geq 0$,
% \[
%     \Lambda_n = \{\gamma \in \Gamma: \left| \widehat{f}(\gamma) \right| \geq \eta_n/2\} \cup \{\gamma \in \Gamma: \left| \widehat{f}(\gamma \circ \phi) \right| \geq \eta_n/2\} \cup \{\gamma \in \Gamma: \left| \widehat{f}(\gamma \circ \psi) \right| \geq \eta_n/2\}.
% \]
% Note that by Plancherel's formula, 
% \[
%     \sum_{\gamma \in \Gamma} \left| \widehat{f}(\gamma)\right|^2 = \lVert f \Vert_2^2 \leq 1.
% \]
% Hence 
% \[
%     \left|\{\gamma \in \Gamma: \left| \widehat{f}(\gamma) \right| \geq \eta_n/2\}\right| \leq 4/\eta_n^2. 
% \]
% A similar argument shows that the cardinality of $\Lambda_n$ is bounded above by a constant that depends only on $\eta_n$ and the indices of $\phi(G), \psi(G)$ in $G$.
% Fix $\delta, \epsilon > 0$. We define sequences $\{\eta_n\}_{n \geq 0}$ and $\{\kappa_n\}_{n \geq 0}$ by $\eta_0 = 1$ \textcolor{red}{In Lemma 2.3, we have $\eta \in (0, 1/2]$} and for $n \geq 0$,
%     \[
%          \kappa_n = \nu(\eta_n)^{-1/2} \text{ and } \eta_{n+1} = \min \left \{\eta_n, \frac{\epsilon^2}{4 \kappa_n^2}, \epsilon \nu \left( \frac{\epsilon}{2 \kappa_n} \right) \right\}
%     \]
% 		where $c$ is the implicit constant in \eqref{eq:c}.
% 	\textcolor{red}{Here are previous $\kappa_n$ and $\nu_n$:
% 	\[
%         \kappa_n = \nu(\eta_n)^{-1/2} \text{ and } \eta_{n+1} = \min \left\{ \eta_n, \frac{\delta^6}{100 K^2_n}, \frac{\delta^3}{5(c+1)} \nu \left( \frac{\delta^3}{10 \kappa_n} \right) \right\},
%     \]
%     }
%     \textcolor{red}{Note that, in the new version, $\epsilon$ replaces $O(\delta^3)$.}	

In view of \cref{lem:kernel_a}, for $n \geq 0$, there is a kernel $K_n: G \to [0, \infty)$ such that
    \[
        \widehat{K}_n \geq 0,  \| \widehat{K_n}\|_1 = \lVert K_n \rVert_{\infty} \leq 1/\nu(\eta_n)
    \]
    and $K_n$ is supported on $B(\Lambda_n; \eta_n) \subseteq D(\eta_n)$. 
%     Note that  $B(\Lambda_n, \eta_n)$ is subset of 
%     \[
%         D(\eta_n) := \{t \in G: d(t) \leq \eta_n \}
%     \]
%     where for $t \in G$,
% 		\[
% 			d(t) := \max \left( \| \widehat{f} - \widehat{f_{t}} \|_\infty, \| \widehat{f} - \widehat{f_{\phi(t)}} \|_\infty, \| \widehat{f} - \widehat{f_{\psi(t)}} \|_\infty  \right).
% 		\] 
%     % \textcolor{red}{I changed $D(\eta_n)$ to $D(\eta_n)$ to avoid confusion with $B(\Lambda_n, \eta_n)$.}
    % \textcolor{red}{$K_n$ will be corresponding to $\Lambda = \{\gamma \in \Gamma: |\hat{f}(\gamma)| \text{ or } |\hat{f}(\gamma \circ \phi)| \text{ or } |\hat{f}(\gamma \circ \psi)| \geq \eta_n/2\}$. Then $K_n$ is supported on $B(\Lambda, \eta_n)$ which is a subset of $D(\eta_n)$.}
    % % \[
    % %     D(\eta_n) = \{t \in G: \left\| \widehat{f} - \widehat{f_{s}} \right\|_{\infty} \leq \eta_n \; \text{ for } s \in \{t, \phi(t), \psi(t)\} \},
    % % \]
    % and $\| \widehat{K_n}\|_1 = \lVert K_n \rVert_{\infty} \leq 1/\nu(\eta_n).$ 
    % % \Hnote{By Lemma 2.1, the set $D(\eta_n)$ is  contained in a Bohr-(x,y) set. By Lemma 2.3, there is a kernel $K_n$ supported on that Bohr set.}
    We define
    \[
        f_n = f * K_n.
    \]
	\noindent \textbf{Claim 1:} 
	\[
	\| \widehat{f} - \widehat{f_n} \|_{\infty} = \sup_{\gamma \in \Gamma} \left| \widehat{f}(\gamma) (1-\widehat{K_n}(\gamma)) \right| \leq \eta_n.
	\]
	Indeed, by construction, $K_n$ is supported on $D(\eta_n)$, and every $t \in D(\eta_n)$ satisfies $\left| \widehat{f}(\gamma) (1 - \gamma(t)) \right| \leq \eta_n$ for all $\gamma \in \Gamma$. Therefore, for all $\gamma \in \Gamma$,
% 	Indeed, from the proof of \cref{lem:translate}, $K_n$ is supported on $B(\Lambda; \eta_n)$ where $\Lambda = \{ \gamma \in \Gamma: \left| \widehat{f}(\gamma) \right| > \frac{\eta_n}{2}\}$. 
% 	If $\gamma \not \in \Lambda$ then clearly 
% 	$ \left| \widehat{f}(\gamma) (1-\widehat{K_n}(\gamma)\right| \leq 2 \left| \widehat{f}(\gamma) \right| \leq \eta_n$. If $\gamma \in \Lambda$ then
	\begin{eqnarray*}
	\left| \widehat{f}(\gamma) \right| \left| 1-\widehat{K_n}(\gamma) \right| 
	&\leq& 
	\left| \widehat{f}(\gamma) \right| \int_{G} K_n(x) \left| 1-\overline{\gamma(x)} \right| \, d\mu(x) \\
	&=& \left| \widehat{f}(\gamma) \right| \int_{D(\eta_n)} K_n(x) \left| 1-\overline{\gamma(x)} \right| \, d\mu(x) \\
	&=& \int_{D(\eta_n)} K_n(x) \left| \widehat{f}(\gamma) \right| \left| 1-\overline{\gamma(x)} \right| \, d\mu(x)\\
	&\leq& \eta_n \int_{D(\eta_n)} K_n(x) \, d\mu(x) \leq \eta_n.
	\end{eqnarray*}
% 	In any case the claim is proved.
	
\noindent	\textbf{Claim 2:} 
    \[
        \lVert f_{n+1} - f_n \rVert_2^2 \leq \lVert f_{n+1} \rVert_2^2 - \lVert f_n \rVert_2^2 + 2 \eta_{n+1} \kappa_n^2.
    \]
	Indeed, we have
	\begin{eqnarray*}
	\lVert f_{n+1} - f_n \rVert_2^2 
&= &	\lVert \widehat{f_{n+1}} - \widehat{f_n} \rVert_2^2 \\
&= &	\lVert \widehat{f_{n+1}} \rVert_2^2 + \lVert \widehat{f_n} \rVert_2^2 - \sum_{\gamma \in \Gamma} 
\left(\widehat{f_{n+1}}(\gamma) \overline{\widehat{f_{n}}(\gamma)} + \overline{\widehat{f_{n+1}}(\gamma)} \widehat{f_{n}}(\gamma) \right)\\
&= &	\lVert \widehat{f_{n+1}} \rVert_2^2 - \lVert \widehat{f_n} \rVert_2^2 + 2 \lVert \widehat{f_n} \rVert_2^2 - \sum_{\gamma \in \Gamma}	\left| \widehat{f}(\gamma) \right|_2^2 \widehat{K_n}(\gamma) \widehat{K_{n+1}}(\gamma)\\
&= & \lVert \widehat{f_{n+1}} \rVert_2^2 - \lVert \widehat{f_n} \rVert_2^2 + 2 \sum_{\gamma \in \Gamma}	\left| \widehat{f}(\gamma) \right|_2^2 \widehat{K_n}(\gamma) 
\left( \widehat{K_n}(\gamma) - \widehat{K_{n+1}}(\gamma) \right)\\
& \leq & \lVert \widehat{f_{n+1}} \rVert_2^2 - \lVert \widehat{f_n} \rVert_2^2 + 2 \sum_{\gamma \in \Gamma}	\left| \widehat{f}(\gamma) \right|^2 \widehat{K_n}(\gamma) 
\left(1 - \widehat{K_{n+1}}(\gamma) \right) \\
& \leq & \lVert \widehat{f_{n+1}} \rVert_2^2 - \lVert \widehat{f_n} \rVert_2^2 + 2 \sup_{\gamma \in \Gamma} |\widehat{f}(\gamma)| \left( 1 - \widehat{K_{n+1}}(\gamma) \right) \cdot \| \widehat{K_n} \|_1 \\
& \leq & \lVert \widehat{f_{n+1}} \rVert_2^2 - \lVert \widehat{f_n} \rVert_2^2 + 2 \eta_{n+1} \kappa_n^2
	\end{eqnarray*}
 and the claim is proved.
   
		Since $\eta_{n+1} \leq \epsilon^2/(4 \kappa_n^2)$, we have
    \[
        \lVert f_{n+1} - f_n \rVert_2^2 \leq \lVert f_{n+1} \rVert_2^2 - \lVert f_n \rVert_2^2 + \epsilon^2/2.
    \]
     Let $M$ be the smallest integer such that $M \geq 2/\epsilon^2$. Then because $\lVert f_M \rVert_2^2 \leq \lVert f_M \rVert_{\infty}^2 \leq 1$,
    \[
        \sum_{n=0}^{M-1} \lVert f_{n+1} - f_n \rVert_2^2 \leq \lVert f_M \rVert_2^2 - \lVert f_0 \rVert_2^2 + M \epsilon^2/2 \leq 1 + M \epsilon^2/2 \leq M \epsilon^2.
    \]
    Therefore there exists $0 \leq n \leq M - 1$ such that
    \[
        \lVert f_{n+1} - f_n \rVert_2 \leq \epsilon.
    \]
    From now on, we fix this $n$.
    Next consider the expression
    \[
        I_n(t) = \int_G f_n(x) f_n(x+\phi(t)) f_n(x + \psi(t)) \, d \mu(x) \; \; \text{ for } t \in G.
    \]
    By the same algebra as in the proof of \cref{lem:J},
    \[
        |I_n(0) - I_n(t)| \leq \| f_n - (f_n)_{\phi(t)} \|_1 + \| f_n - (f_n)_{\psi(t)} \|_1. 			
    \]
	Note that
	\begin{eqnarray*}
		\| f_n - (f_n)_{\phi(t)} \|^2_1 &\leq & \| f_n - (f_n)_{\phi(t)} \|_2^2 = \| (f - f_{\phi(t)})*K_n \|_2^2 \\
		&=& \sum_{\gamma \in \Gamma} \left| \widehat{K_n}(\gamma) \right|^2 \left|\widehat{f}(\gamma) - \widehat{f_{\phi(t)}}(\gamma) \right|^2 \\
		&\leq & \| \widehat{K_n} \|_1 d(t)^2 \leq \kappa_n^2 d(t)^2.
		\end{eqnarray*}
		The same estimate holds for $\| f_n - (f_n)_{\psi(t)} \|_1^2$. Hence $|I_n(0) - I_n(t)| \leq 2 \kappa_n d(t)$ for any $t \in G$.
		
    Since $I_n(0) = \lVert f_n^3 \rVert_1 \geq \lVert f_n \rVert_1^3 = \lVert f \rVert_1^3 \geq \delta^3$, it follows that
    \[
        I_n(t) \geq \delta^3 - 2 \kappa_n d(t) \text{ for all } t \in G.
    \]
    Note that $d(t) \leq \epsilon/(2\kappa_n)$ for $t$ in the set $D(\epsilon/(2 \kappa_n))$ and so $I_n(t) \geq \delta^3 - \epsilon$ in this set. 
    
    %Since $\eta_{n+1} \leq \epsilon \nu \left( \frac{\epsilon}{2 \kappa_n} \right)$, we can choose a constant $\eta < \frac{\epsilon}{2 \kappa_n}$ such that $\eta_{n+1} < \epsilon \nu(\eta)$. Let $k$ be the kernel corresponding to this $\eta$ found in \cref{lem:kernel_a}. Then $\lVert k \rVert_{\infty} \leq 1/\nu(\eta)$ and $k$ is supported on $B(\eta) \subseteq B(\epsilon/(2\kappa_n))$. 
    
    Let $\eta = \epsilon/(2 \kappa_n)$. In view of \cref{lem:kernel_a}, there exists a kernel $K$ supported on $D(\eta)$ such that $\lVert K \rVert_{\infty} \leq 1/\nu(\eta)$.
    % Pick $\eta = \frac{\epsilon}{2 \kappa_n}$ and let $K$ be the kernel corresponding to this $\eta$ and being supported on $B(\eta)$. The existence of $K$ is guaranteed by \cref{lem:kernel_a}. Then $\lVert K \rVert_{\infty} \leq 1/\nu(\eta)$ and $K$ is supported on $B(\eta) = B(\epsilon/(2\kappa_n))$. 
    We then have
    \[
        J'(f_{n}) := \int_G I_{n}(t) K(t) \, d \mu(t) \geq (\delta^3 - \epsilon) \int_{D(\eta)} K(t) \, d\mu(t) \geq \delta^3 - \epsilon.
    \]
    Letting $f_{st} = f_n$, $f_{er} = f_{n+1} - f_n$ and $f_{un} = f - f_{n+1}$, we obtain
    \begin{enumerate}
        \item $\lVert K \rVert_{\infty} \leq 1/\nu(\epsilon/(2\kappa_n)) \leq 1/\nu(\epsilon/(2 \kappa_M))$ (since $\nu$ is increasing on $(0, e^{1/2}/C_0) \supset (0, 1)$),

        \item $\lVert f_{st} \rVert_{\infty} = \lVert f_n \rVert_{\infty} = \lVert f * K_n \rVert_{\infty} \leq \lVert f \rVert_{\infty} \lVert K_n \rVert_1 \leq 1$, and $\lVert f_{er} \rVert_{\infty} \leq 2$, and $\lVert f_{un} \rVert_{\infty} \leq 2$,
        
        \item $\lVert f_{er} \rVert_2 = \lVert f_{n+1} - f_n \rVert_2 \leq \epsilon$,
        
        \item $\lVert \hat{f}_{un} \rVert_{\infty} \lVert K \rVert_{\infty} = \lVert \hat{f} - \hat{f}_{n+1} \rVert_{\infty} \lVert K \rVert_{\infty} < \eta_{n+1}/ \nu(\eta) \leq \epsilon$ because $\eta_{n+1} \leq \epsilon \nu(\epsilon/(2\kappa_n)) = \epsilon \nu(\eta).$
        
        \item $J'(f_{st}) = J'(f_n) \geq \delta^3 - \epsilon$.
    \end{enumerate}
    Our proof finishes. 
\end{proof}

\subsection{Proof of density results}
The goal of this section is to prove \cref{th:main-density} and \cref{th:roth-khintchine}. First we recall \cref{th:roth-khintchine} for the reader's convenience.

\begin{theorem*}[Khintchine-Roth theorem for compact abelian groups]  
Let $f: G \to [0, 1]$ be a measurable function with $\int_G f \, d \mu > \delta$. Let $\phi, \psi: G \to G$ be continuous homomorphisms such that $[G: \phi(G)], [G: \psi(G)]$ and $[G:(\phi-\psi)G]$ are finite. Then for every $\epsilon > 0$, there exists a constant $c_1$ that depends only on $\delta, \epsilon$ and the indexes above such that the set
\[
    B = \left\{t \in G: \int_G f(x) f(x+ \phi(t)) f(x + \psi(t)) \, d\mu(x) > \delta^3 - \epsilon \right\}
\]
has measure greater than $c_1$. As a consequence, there exists a constant $c_2$ that depends only on $\delta$ and indexes of $\phi(G), \psi(G), (\phi - \psi)(G)$ such that
\[
    J(f) := \iint_{G^2} f(x) f(x+ \phi(t)) f(x+ \psi(t)) d\mu(x) d\mu(t) > c_2.
\]
\end{theorem*}
\begin{proof}
    Fix $\epsilon > 0$ and let constant $C$, kernel $K$ and the decomposition $f = f_{st} + f_{er} + f_{un}$ be as found in \cref{prop:regularity_lemma}. Define
    \[
        J'(f) := \iint_{G^2} f(x) f(x+ \phi(t)) f(x+ \psi(t)) K(t) \ d\mu(x) d\mu(t)
    \]
    and 
    \[
        J'(f_{st}) := \iint_{G^2} f_{st}(x) f_{st}(x + \phi(t)) f_{st}(x + \psi(t)) K(t)\ d\mu(x) d\mu(t).
    \]
    
    Applying the decomposition $f = f_{st} + f_{er} + f_{un}$ and expanding $J'(f)$, we see that the difference $J'(f) - J'(f_{st})$ will have $26$ terms. The terms that contain $f_{er}$ can be bounded by $4 \epsilon$ since for $f_1, f_2, f_3 \in L^{\infty}(G)$,
    \begin{equation*}
        \iint_{G^2} f_1(x) f_2(x + \phi(t)) f_3(x + \psi(t)) K(t) \, d\mu(x) d\mu(t) \leq \max_{i} \lVert f_i \rVert_{\infty}^2 \max_{j} \lVert f_j \rVert_1 \lVert K \rVert_1.
    \end{equation*}
    
    On the other hand, in view of \cref{lem:bourgain_lemma_2} and a change of variables if necessary, the terms containing $f_{un}$ are bounded by $O(\lVert \hat{f}_{un} \rVert_{\infty} \lVert \widehat{K} \rVert_{1})$ which is $O(\epsilon)$ thanks to the properties of the decomposition. (The implicit constant may depend on the index $[G:(\phi - \psi)(G)]$ now due to this change of variables.)
    %Note that the implicit constant depends only on the indices of $\phi(G)$ and $\psi(G)$. 
    Therefore,
    \begin{equation}
    \label{eq:J_f'_1}
        J'(f) > J'(f_{st}) - O(\epsilon) > \delta^3 - c_0 \epsilon
    \end{equation}
    where the constant $c_0$ depends only on the indexes of $\phi(G), \psi(G)$ and $(\phi - \psi)(G)$ in $G$. 
    
    Define
    \begin{equation*}
        I_f(t) = \int_G f(x) f(x+ \phi(t)) f(x+ \psi(t)) \, d\mu(x)
    \end{equation*}
    and 
    \begin{equation*}
        B = \{t \in G: I_f(t) > \delta^3 - 2 c_0 \epsilon \}. 
    \end{equation*}
    We then have
    \begin{align}
    \label{eq:J_f'_2}
    \begin{aligned}
        J'(f) = \int_G I_f(t) K(t) \, d\mu(t) = \int_B I_f(t) K(t) \, d\mu(t) + \int_{G \setminus B} I_f(t) K(t) \, d\mu(t) \leq \\ \int_B K(t) \, d\mu(t) + (\delta^3 - 2 c_0 \epsilon) \int_{G \setminus B} K(t) \, d\mu(t) \leq 
        \lVert K \rVert_{\infty} \mu(B) + (\delta^3 - 2 c_0 \epsilon).
    \end{aligned}
    \end{align}
    Combining \eqref{eq:J_f'_1} and \eqref{eq:J_f'_2}, we deduce that 
    \begin{equation*}
        \mu(B) > c_0 \epsilon/\lVert K \rVert_{\infty}.
    \end{equation*}
    Letting $c_1 = c_0 \epsilon/\lVert K \rVert_{\infty}$, we obtain the first part of the theorem. 
    
    Now we have
    \begin{equation*}
        J(f) = \int_G I_f(t) \, d\mu(t) > (\delta^3 - 2 c_0 \epsilon) c_1.
    \end{equation*}
    Letting $c_2 = c_1(\delta^3 - 2c_0 \epsilon)$, we obtain the second part of the theorem. 
\end{proof}


In order to prove \cref{th:main-density}, we need the following proposition. For our future applications, we will state and prove a slightly more general version than what is necessary.

% For our future applications, it is convenient to state the following fact separately and slightly stronger than what we need:

\begin{proposition}
\label{prop:density-bohr-stronger}
Suppose $\phi, \psi: G \rightarrow G$ are continuous homomorphisms such that $\phi(G), \psi(G), (\phi-\psi)(G)$ have finite index in $G$. Let $f : G \rightarrow [0,1]$ such that $\int_G f \, d\mu = \delta >0$. For $w \in G$, define
\[
    R(w) = \iint_{G^2} f(x + w) f(x + \phi(y)) f(x+\psi(y))  \, d\mu(x) d\mu(y)
\]
Then there are $c, k, \eta >0$ depending only on $\delta$ and the indexes above such that the set
 $\{ w \in G: R(w) > c\}$ contains a Bohr-$(k,\eta)$ set. 
\end{proposition}


\begin{proof}
By \cref{lem:bourgain_lemma_2}, we have
\[
    |R(w) - R(0)| \ll \lVert \hat{f} - \widehat{f_w} \rVert_{\infty}
\]
where implicit constant depends only on the indexes of $\phi(G), \psi(G)$ in $G$. 
By \cref{th:roth-khintchine}, we know that $R(0) > c$ for some constant $c>0$ depending on the indexes $[G: \phi(G)]$, $[G: \psi(G)]$, $[G: (\phi - \psi)(G)]$ and $\delta$. It follows that there exists a constant $c'$ such that $R(w) > c/2$ if 
\begin{equation}
% \label{eq:br_f_w_needed}
    \lVert \hat{f} - \widehat{f_w} \rVert_{\infty} < c'.
\end{equation}
\cref{lem:translate} implies that the set of such $w$ contains a Bohr-$(k,\eta)$ set, where $k$ and $\eta$ depend only on $c'$.
\end{proof}

We can now formulate and prove our main theorem for sets of positive measure.

\begin{theorem} \label{th:main-density-true}
Let $\phi_1, \phi_2, \phi_3, \psi_1, \psi_2: G \rightarrow G$ be continuous homomorphisms satisfying the following:
\begin{enumerate}
\item $\phi_1 + \phi_2 + \phi_3=0$,
\item $\phi_1 \circ \psi_1 = \phi_2 \circ \psi_2$,
\item $\phi_3(G), \psi_1(G), \psi_2(G), (\psi_1 + \psi_2)(G)$ have finite index in $G$.
\end{enumerate}
The for any measurable set $A \subset G$, $\mu(A) = \delta >0$, the set $\phi_1(A) + \phi_2(A) + \phi_3(A)$ contains a Bohr-$(k, \eta)$ set, where $k$ and $\eta$ depend only on $\delta$ and the indexes above. 
\end{theorem}
\begin{proof}
Applying \cref{prop:density-bohr-stronger} for $f=1_A$ and $\psi_1, -\psi_2$ in place of $\phi$ and $\psi$, we see that there exists a Bohr-$(k,\eta)$ set $B$ such that for all $w \in B$, there are $x, y \in G$ such that
\[
x+w , x+\psi_1(y) \textup{ and } x-\psi_{2}(y) \in A.
\]
Note that
\[
\phi_3(x + w) + \phi_1(x+\psi_1(y)) + \phi_2(x-\psi_2(y)) = \phi_3(w)
\]
and so that $\phi_1(A) + \phi_2(A) + \phi_3(A) \supseteq \phi_3(B)$.
Our theorem then follows from \cref{lem:bohr-homo2}.
\end{proof}

\cref{th:main-density} is now a special case of \cref{th:main-density-true} with $\psi_1 = \phi_2$ and $\psi_2 = \phi_1$. Note that the condition of \cref{th:main-density-true} is met because $\phi_1 + \phi_2 = - \phi_3$ and so
\[
    [G: (\psi_1 + \psi_2)(G)] = [G:(\phi_1 + \phi_2)(G)] = [G:(-\phi_3)(G)] = [G: \phi_3(G)]
\]
which is finite.

